% =====================================================================
%  패키지와 서식 설정
% =====================================================================

% ----- 한글 -----
\usepackage[hangul]{kotex}      % 한국어 조판: "제 1 장", "그림", "표", "목차" 등
% 글꼴을 바꾸려면 아래 주석을 풀고 설치된 글꼴 이름을 넣으세요
% \setmainhangulfont{NanumMyeongjo}
% \setsanshangulfont{NanumGothic}

% ----- 쪽 여백 (학교 규정에 맞게 수정) -----
\usepackage[top=30mm,bottom=30mm,left=30mm,right=30mm]{geometry}

% ----- 줄 간격 -----
\usepackage{setspace}
\setstretch{1.6}

% ----- 수식 -----
\usepackage{amsmath,amssymb,amsthm}

% ----- 그림·표 -----
\usepackage{graphicx}
\graphicspath{{../figures/}}    % figures/ 폴더의 그림을 파일 이름만으로 불러옴
\usepackage{booktabs}           % 깔끔한 표 선 (\toprule, \midrule, \bottomrule)
\usepackage{caption}
\captionsetup{labelsep=period}
\usepackage{float}

% ----- 참고문헌 (스타일은 학교·학과 규정에 맞게 변경) -----
\usepackage{csquotes}
\usepackage[backend=biber,style=authoryear,sorting=nyt]{biblatex}
\addbibresource{../references/references.bib}

% ----- 링크·상호 참조 (가장 마지막에 불러옴) -----
\usepackage[hidelinks]{hyperref}
